\chapter{Introduction}
\label{ch:introduction}

Can carefully engineered classical chess search execute efficiently inside a managed JavaScript runtime? This is the central question that motivated the Kronos project. JavaScript's ubiquity in modern software development makes it an attractive target for real-time, browser-deployable research tools. Yet the Google V8 managed runtime introduces two engineering constraints that have historically made it an unlikely candidate for performance-critical game-tree search~\cite{richards2010analysis, selakovic2016empirical}.

The first constraint is the \textbf{single-threaded event loop}. In a browser or Node.js environment, all computation—UI rendering, network handlers, and game logic—shares one thread. Running a minimax search on this thread causes the interface to freeze for the duration of the search, making responsive, interactive play impossible. The second constraint is \textbf{garbage collection (GC) latency}. Chess search expands millions of nodes per second, each generating short-lived move objects, evaluation records, and hash map entries. Without careful memory discipline, the V8 GC triggers major collection sweeps that introduce latency spikes of hundreds of milliseconds, degrading nodes-per-second (NPS) throughput unpredictably~\cite{jones2011garbage}.

To address both constraints simultaneously, we built Kronos: a modular chess engine and automated benchmarking and validation suite implemented in Node.js and JavaScript. Kronos decouples the search core into background Web Workers, communicating with the React UI via serialized message passing. Inside the search path, we eliminate transient object allocations through in-place board mutations, flat TypedArray move buffers, and a size-bounded transposition table, keeping the V8 heap stable between 32.85\,MB and 445.66\,MB during tournament play (with single-game telemetry profiling showing active heap usage down to 24.98\,MB), building on concurrency patterns for dynamic runtimes~\cite{hunter2021multithreaded}.




\begin{figure}[htbp]
\centering
\begin{tikzpicture}[
    node distance=1.2cm and 1.8cm,
    block/.style={draw=blue!60!black, rectangle, minimum height=3em, minimum width=10em,
                  text width=9.5em, align=center, rounded corners=3pt, fill=blue!10, font=\small\sffamily},
    arrow/.style={draw, -latex, thick, gray!80!black}
]
    \node [block, fill=gray!15, draw=gray!60!black] (p1) {Phase 1\\Minimax / Alpha-Beta};
    \node [block, fill=blue!10, draw=blue!60!black, right=of p1] (p2) {Phase 2\\Move Ordering \& PSTs};
    \node [block, fill=blue!10, draw=blue!60!black, right=of p2] (p3) {Phase 3\\LMR, NMP \& Memory};
    \node [block, fill=blue!10, draw=blue!60!black, below=1.2cm of p3] (p4) {Phase 4\\SPRT Validation};

    \draw [arrow] (p1) -- (p2);
    \draw [arrow] (p2) -- (p3);
    \draw [arrow] (p3) -- (p4);

    % Elo labels below each phase
    \node[font=\scriptsize\sffamily, text=gray!80!black, below=0.1cm of p1] {$\sim$900 Elo};
    \node[font=\scriptsize\sffamily, text=gray!80!black, below=0.1cm of p2] {$\sim$1{,}380 Elo};
    \node[font=\scriptsize\sffamily, text=gray!80!black, below=0.1cm of p3] {$\sim$1{,}468 Elo};
    \node[font=\scriptsize\sffamily, text=blue!80!black, below=0.1cm of p4] {\textbf{1{,}485 Elo}};

\end{tikzpicture}
\caption[Development Timeline and Optimization Phases]{Chronological Development Timeline and Optimization Phases of the Kronos Engine Project.}
\label{fig:engine_timeline}
\end{figure}

% Source Material: ROADMAP_NEXT.md and README.md


Figure~\ref{fig:engine_timeline} illustrates the chronological phases of our development, tracing the progression from a basic Minimax implementation to the fully optimized configuration. We focus on empirical measurement under controlled tournament conditions rather than theoretical performance bounds. By tracking execution metrics across thousands of simulated matches, we isolate the specific search space compression and relative playing strength shifts associated with each heuristic. These measurements identify which engineering optimizations are most effective under the V8 memory manager and establish Kronos as a reproducible platform for dynamic language systems research.

The remainder of this report is structured as follows. Chapter~\ref{ch:contributions} details our core research questions and engineering contributions. Chapter~\ref{ch:related_work} surveys game-tree search background and compares Kronos with other JavaScript chess engines. Chapter~\ref{ch:system_overview} details the architecture, memory model, and worker thread design. Chapter~\ref{ch:experimental_methodology} outlines the tournament benchmarking pipeline and statistical model. Chapter~\ref{ch:experimental_results} presents our empirical results, while Chapter~\ref{ch:discussion} analyzes the outcomes and explores self-play draw rates. Chapter~\ref{ch:limitations} reviews threats to validity, Chapter~\ref{ch:future_work} suggests next steps, and Chapter~\ref{ch:conclusion} summarizes the broader contribution.


